\documentclass[a4paper]{article}
\usepackage{gpti}

% Tabelas e números gerados por ahp.py em tabelas/: não edite à mão.
% Sem eles, o documento compila com marcas de pendência no lugar.
\newcommand{\pendente}{\textcolor{red}{[PENDENTE]}}
\IfFileExists{tabelas/numeros.tex}{\newcommand{\crCriterios}{0,002}
\newcommand{\crCusto}{0,000}
\newcommand{\crPrazo}{0,027}
\newcommand{\crManutencao}{0,004}
\newcommand{\crMaoDeObra}{0,005}
\newcommand{\pesoCusto}{7,0\%}
\newcommand{\pesoPrazo}{36,8\%}
\newcommand{\pesoManutencao}{19,3\%}
\newcommand{\pesoMaoDeObra}{36,8\%}
\newcommand{\primeiroLugar}{Java}
\newcommand{\pontuacaoPrimeiro}{34,3\%}
\newcommand{\segundoLugar}{Python}
\newcommand{\pontuacaoSegundo}{33,6\%}
}{
    \newcommand{\crCriterios}{\pendente}
    \newcommand{\crCusto}{\pendente}
    \newcommand{\crPrazo}{\pendente}
    \newcommand{\crManutencao}{\pendente}
    \newcommand{\crMaoDeObra}{\pendente}
    \newcommand{\pesoCusto}{\pendente}
    \newcommand{\pesoPrazo}{\pendente}
    \newcommand{\pesoManutencao}{\pendente}
    \newcommand{\pesoMaoDeObra}{\pendente}
    \newcommand{\primeiroLugar}{\pendente}
    \newcommand{\pontuacaoPrimeiro}{\pendente}
    \newcommand{\segundoLugar}{\pendente}
    \newcommand{\pontuacaoSegundo}{\pendente}
}
% Cabeçalhos das tabelas geradas quebram entre palavras, nunca no meio de uma
\newcommand{\nohyphens}{\hyphenpenalty=10000\exhyphenpenalty=10000}
\newcommand{\tabelagerada}[1]{\IfFileExists{tabelas/#1.tex}{\input{tabelas/#1}}{\textcolor{red}{[PENDENTE: rodar ahp.py]}}}

\hyphenation{Know-ledge}

\title{\textbf{Grupo 7 - AHP: escolha da ferramenta do trabalho semestral}}

\begin{document}

\capa

\section{Objetivo e hierarquia}

O objetivo é escolher, com o AHP (Analytic Hierarchy Process), a ferramenta para implementar o trabalho semestral: o Sistema de Controle de Acesso da EACH, que controla a entrada e a saída de pessoas e veículos, libera automaticamente a comunidade pelo Nº USP ou CPF, cadastra visitantes, tem hierarquia de acesso por senha, faz o pré-cadastro em lote para grandes eventos, emite relatórios de frequência e se integra às catracas e cancelas. A escolha é o trabalho da atividade K, \enquote{Definir a stack tecnológica} (EAP 1.2.2.1, no pacote Documento de arquitetura), do cronograma do grupo (Exercício 2).

O AHP organiza a decisão em uma hierarquia, com o objetivo no topo, os critérios no meio e as alternativas na base (Figura~\ref{fig:hierarquia}) \cite{saaty, vargas2010}. É um caso de análise multicritério ponderada, que o PMBOK cita como exemplo de esquema de priorização \cite[p.~181]{pmi2021}.

\begin{figure}[H]
    \centering
    \begin{tikzpicture}[
            objetivo/.style={fill=eapA, text=white, font=\sffamily\small\bfseries, inner sep=5pt, rounded corners=1pt},
            criterio/.style={fill=eapB, text=white, font=\sffamily\scriptsize\bfseries, text width=3cm,
                             minimum height=9mm, align=center, inner sep=3pt, rounded corners=1pt},
            alternativa/.style={draw=eapA, fill=eapC, font=\sffamily\scriptsize, text width=3cm,
                                minimum height=8mm, align=center, inner sep=3pt, rounded corners=1pt},
            nivel/.style={font=\sffamily\scriptsize\itshape, text=eapA, anchor=east},
            liga/.style={draw=eapA, semithick}]
        \foreach \nome [count=\i] in {Custo de desenvolvimento, Prazo de entrega, Facilidade de manutenção,
                                      Facilidade de mão de obra}
            \node[criterio] (c\i) at ({3.8*(\i-1)}, 0) {\nome};
        \foreach \nome [count=\i] in {Java / Spring Boot, C\# / ASP.NET Core, Python / Django, Go / Gin}
            \node[alternativa] (a\i) at ({3.8*(\i-1)}, -2.3) {\nome};
        \node[objetivo] (obj) at (5.7, 1.6) {Escolher a ferramenta do Sistema de Controle de Acesso da EACH};
        \foreach \i in {1,...,4}
            \draw[liga] (obj.south) -- ++(0, -0.4) -| (c\i.north);
        % Uma faixa no lugar das 16 ligações: todas as alternativas passam por todos os critérios
        \node[draw=eapA, dashed, rounded corners=1pt, minimum width=14.6cm, minimum height=5mm,
              font=\sffamily\scriptsize\itshape, text=eapA] (faixa) at (5.7, -1.15)
            {cada alternativa é comparada com as outras sob cada um dos critérios};
        \foreach \i in {1,...,4} {
            \draw[liga] (c\i.south) -- (c\i.south |- faixa.north);
            \draw[liga] (a\i.north |- faixa.south) -- (a\i.north);
        }
        \node[nivel] at ($(c1.west |- obj)+(-0.3, 0)$) {Objetivo};
        \node[nivel] at ($(c1.west)+(-0.3, 0)$) {Critérios};
        \node[nivel] at ($(a1.west)+(-0.3, 0)$) {Alternativas};
    \end{tikzpicture}
    \caption{Hierarquia da decisão.}
    \label{fig:hierarquia}
\end{figure}

\section{Premissas e critérios}

\begin{itemize}
    \item A equipe é a de bolsistas: 5 a 8 alunos do último ano de SI, com o Analista de Sistemas Sênior. Os Exercícios 1 e 2 deixaram em aberto a escolha entre bolsistas e uma start-up, e o resultado vale para a equipe de bolsistas.
    \item Cada alternativa é a linguagem com o seu principal framework web: Java/Spring Boot, C\#/ASP.NET Core, Python/Django e Go/Gin.
\end{itemize}

A Tabela~\ref{tab:definicoes} define os critérios. Cada fator entra em um único critério, para não ser contado duas vezes: a curva de aprendizado, por exemplo, só pesa na mão de obra, porque o prazo já supõe um time que sabe a linguagem.

\begin{table}[H]
    \centering
    \small
    \renewcommand{\arraystretch}{1.3}
    \begin{tabular}{|>{\raggedright\arraybackslash}p{2.8cm}|>{\raggedright\arraybackslash}p{9.6cm}|>{\raggedright\arraybackslash}p{3.8cm}|}
        \hline
        \rowcolor{eapC}
        \textbf{Critério} & \textbf{O que mede} & \textbf{Base} \\
        \hline
        \textbf{Custo de desenvolvimento} & Gastos em R\$ que mudam com a linguagem: licenças, ferramentas, hospedagem e integração com catracas e cancelas. As bolsas são fixas e ficam de fora. & PMBOK 6ª ed., cap. 7 \cite{pmi2017} \\
        \hline
        \textbf{Prazo de entrega} & O que o framework traz pronto para entregar os quatro módulos (controle de acesso, gestão de usuários e visitantes, rotinas especiais e relatórios), supondo um time que já sabe a linguagem. & PMBOK 6ª ed., cap. 6 \cite{pmi2017} \\
        \hline
        \textbf{Facilidade de manutenção} & Depois da implantação (a manutenção ficou fora do escopo do projeto, Exercício 1), quão fácil é para quem herdar o código corrigir, evoluir e adaptar o sistema. A norma define a manutenibilidade como \enquote{Capability of a product to be modified by the intended maintainers with effectiveness and efficiency}. O tempo de suporte das versões é um acréscimo deste trabalho, não da norma. & ISO/IEC 25010:2023 \cite{arc42iso} \\
        \hline
        \textbf{Facilidade de mão de obra} & Durante o projeto, achar bolsistas de SI que já sabem a linguagem, a curva de aprendizado de quem não sabe e a reposição de quem sair. & PMBOK 6ª ed., 9.3 Adquirir Recursos: critérios de seleção como disponibilidade, experiência e conhecimento \cite{pmi2017} \\
        \hline
    \end{tabular}
    \caption{Definição dos critérios.}
    \label{tab:definicoes}
\end{table}

\section{Método}

Em cada matriz, o elemento da linha é comparado ao da coluna pela escala fundamental de Saaty (Tabela~\ref{tab:escala}) \cite{saaty}. A comparação inversa recebe o recíproco: se Java vale 3 diante de Go, Go vale 1/3 diante de Java.

\begin{table}[H]
    \centering
    \small
    \renewcommand{\arraystretch}{1.25}
    \begin{tabular}{|>{\centering\arraybackslash}p{3.4cm}|>{\raggedright\arraybackslash}p{9cm}|}
        \hline
        \rowcolor{eapC}
        \textbf{Valor} & \textbf{Significado} \\
        \hline
        1 & Igual importância \\
        \hline
        3 & Importância moderada \\
        \hline
        5 & Importância essencial ou forte \\
        \hline
        7 & Importância muito forte \\
        \hline
        9 & Importância extrema \\
        \hline
        2, 4, 6, 8 & Valores intermediários entre dois julgamentos vizinhos \\
        \hline
        1/2 a 1/9 & Recíprocos, para a comparação inversa \\
        \hline
    \end{tabular}
    \caption{Escala fundamental de Saaty.}
    \label{tab:escala}
\end{table}

O vetor de prioridades de cada matriz é o seu autovetor principal \cite{saaty}. A ideia do cálculo aparece na aproximação usada por Vargas \cite{vargas2010}: dividir cada valor pela soma da sua coluna e tirar a média de cada linha. Essa média aproxima a prioridade de cada elemento; as tabelas trazem o autovetor exato, calculado por programa, e as prioridades de cada matriz somam 100\%.

Os julgamentos podem se contradizer: se Java vence C\# e C\# vence Go, Go não deveria vencer Java. A consistência é medida pelo autovalor principal $\lambda_{\max}$, que é igual a $n$, o número de elementos comparados, quando a matriz é perfeitamente consistente \cite{saaty}:
\[
    \mathrm{CI} = \frac{\lambda_{\max} - n}{n - 1}
    \qquad\qquad
    \mathrm{CR} = \frac{\mathrm{CI}}{\mathrm{RI}}
\]
O CI é o índice de consistência, e o CR (Consistency Ratio) o compara ao índice aleatório RI, que vale 0,90 para $n = 4$. A matriz é consistente quando CR $<$ 0,10 \cite{vargas2010}.

A pontuação global de cada alternativa $a$ soma, sobre os critérios $c$, o peso do critério vezes a prioridade da alternativa naquele critério \cite{vargas2010}:
\[
    G(a) = \sum_{c} w_c \, p_c(a)
\]
em que $w_c$ é o peso do critério $c$ e $p_c(a)$ é a prioridade de $a$ sob $c$.

\section{Pesos dos critérios}

A Tabela~\ref{tab:criterios} compara os critérios dois a dois quanto à importância para o objetivo. Os pesos ficam em \pesoCusto{} para custo de desenvolvimento, \pesoPrazo{} para prazo de entrega, \pesoManutencao{} para facilidade de manutenção e \pesoMaoDeObra{} para facilidade de mão de obra, com CR de \crCriterios.

Prazo de entrega e facilidade de mão de obra empatam como os critérios mais pesados, porque a equipe só existe enquanto duram as bolsas. No cronograma do Exercício 2, as bolsas cobrem cerca de 26 semanas a partir da formalização da equipe, da semana 10 à 35; o projeto termina na semana 26, e as 9 semanas restantes são a única reserva para atrasos. Uma ferramenta que atrase a entrega consome essa reserva, e uma equipe que não conheça a linguagem gasta parte dessas mesmas semanas aprendendo. A facilidade de manutenção vem em seguida: a manutenção ficou fora do escopo, mas o sistema continua em uso depois da implantação. O custo pesa pouco, porque as quatro linguagens e seus frameworks têm licenças gratuitas e as bolsas são fixas.

\begin{table}[H]
    \centering
    \small
    \renewcommand{\arraystretch}{1.25}
    \tabelagerada{matriz_criterios}
    \caption{Matriz de comparação dos critérios e vetor de pesos (CR~=~\crCriterios).}
    \label{tab:criterios}
\end{table}

\section{Alternativas em cada critério}

Em cada critério, as quatro alternativas são comparadas duas a duas, sempre para a equipe de bolsistas descrita nas premissas. O PMBOK põe a escolha da ferramenta entre as decisões de grande impacto de um projeto e lembra que decidir em grupo aproveita o conhecimento de mais pessoas e aumenta o compromisso com o resultado \cite[p.~27--28]{pmi2021}. Por isso cada matriz vem com o motivo dos seus julgamentos e a fonte de cada fato, para que qualquer membro possa conferir e contestar um número.

\subsection{Custo de desenvolvimento}

\begin{table}[H]
    \centering
    \small
    \renewcommand{\arraystretch}{1.25}
    \tabelagerada{matriz_custo}
    \caption{Matriz de comparação e vetor de prioridades sob o critério custo de desenvolvimento (CR~=~\crCusto).}
    \label{tab:custo}
\end{table}

A Tabela~\ref{tab:custo} tem CR de \crCusto. As quatro linguagens e seus frameworks são de código aberto. A diferença está nas ferramentas do C\#. Em organizações com mais de 250 computadores ou mais de US\$~1 milhão de receita anual, o Visual Studio Community e o C\# Dev Kit só podem ser usados em projetos de código aberto, pesquisa acadêmica e sala de aula \cite{vscommunity, csdevkit}. A USP passa desse porte, com receita prevista de R\$~9,41 bilhões em 2026 \cite{uspdiretrizes}, e um sistema de controle de acesso não se encaixa claramente em nenhum desses casos. O caminho gratuito é a extensão C\# do VS Code, que pode ser usada sozinha \cite{csharpext}, mas sem a visão de projeto e o explorador de testes que o Dev Kit acrescenta; o Dev Kit completo vem com as assinaturas pagas do Visual Studio Professional e Enterprise \cite{csdevkit}. Por isso o C\# leva uma pequena desvantagem (2) diante das outras três, que empatam entre si.

\subsection{Prazo de entrega}

\begin{table}[H]
    \centering
    \small
    \renewcommand{\arraystretch}{1.25}
    \tabelagerada{matriz_prazo}
    \caption{Matriz de comparação e vetor de prioridades sob o critério prazo de entrega (CR~=~\crPrazo).}
    \label{tab:prazo}
\end{table}

A Tabela~\ref{tab:prazo} tem CR de \crPrazo. O Django gera um painel administrativo a partir dos modelos, pensado para uso interno, e traz usuários, permissões, grupos e telas de login \cite{djangoadmin, djangoauth}: isso cobre boa parte dos cadastros e da hierarquia de acesso. Spring Boot e ASP.NET Core trazem autenticação e controle de acesso por papéis \cite{springsecurity, aspnetidentity} e o acesso ao banco pelo Spring Data JPA e pelo EF Core, com migrações do esquema \cite{springdatajpa, springbootdb, efcoremigrations}, mas não o painel, e empatam entre si. A documentação do Gin trata de rotas, middleware, validação e renderização \cite{gin}, e ORM, migrações e o cadastro de usuários e papéis ficariam com bibliotecas à parte ou com o time.

\subsection{Facilidade de manutenção}

\begin{table}[H]
    \centering
    \small
    \renewcommand{\arraystretch}{1.25}
    \tabelagerada{matriz_manutencao}
    \caption{Matriz de comparação e vetor de prioridades sob o critério facilidade de manutenção (CR~=~\crManutencao).}
    \label{tab:manutencao}
\end{table}

A Tabela~\ref{tab:manutencao} tem CR de \crManutencao. Dois fatores separam as alternativas. O primeiro é a tipagem: Java, C\# e Go têm tipagem estática, verificada pelo compilador \cite{jls, csharptypes, gospec}, e no Python as anotações de tipo são opcionais e não são verificadas na execução \cite{pythontyping}, o que torna mais arriscadas as mudanças de quem herdar o código. O segundo é o tempo de suporte das versões: o Java tem versões LTS gratuitas com pelo menos quatro anos de correções \cite{adoptium}, o .NET e o Django têm versões LTS de três anos \cite{dotnetsupport, djangodownload}, e o Python mantém cada versão por cinco \cite{pythonversions}. Já cada versão do Go só recebe correções até saírem duas mais novas \cite{gorelease}, o que obriga a atualizar com mais frequência, ainda que a promessa de compatibilidade do Go~1 torne essa troca mais simples \cite{go1compat}. Java e C\# vão bem nos dois fatores e empatam. Entre Go e Python, a tipagem pesou mais que o suporte.

\subsection{Facilidade de mão de obra}

\begin{table}[H]
    \centering
    \small
    \renewcommand{\arraystretch}{1.25}
    \tabelagerada{matriz_maodeobra}
    \caption{Matriz de comparação e vetor de prioridades sob o critério facilidade de mão de obra (CR~=~\crMaoDeObra).}
    \label{tab:maodeobra}
\end{table}

A Tabela~\ref{tab:maodeobra} tem CR de \crMaoDeObra. A ACH2003 (Computação Orientada a Objetos), cuja bibliografia é de Java \cite{ach2003}, é disciplina obrigatória do 3º semestre na grade de SI \cite{gradesi}: bolsistas do último ano já passaram por ela. Fora da EACH, entre desenvolvedores profissionais, a pesquisa do Stack Overflow de 2025 põe Python à frente das outras três alternativas, com 54,8\%, contra 29,9\% do C\#, 29,6\% do Java e 17,4\% do Go \cite{so2025}. A pesquisa mede o mercado mundial, não os alunos da EACH.

\section{Ranking global}

Na Tabela~\ref{tab:ranking}, cada célula de critério é a prioridade da alternativa naquele critério, e o total é a soma ponderada pelos pesos do cabeçalho. Os percentuais têm uma casa decimal e podem não somar exatamente 100\%. A ferramenta escolhida é \primeiroLugar, com pontuação global de \pontuacaoPrimeiro, à frente de \segundoLugar, com \pontuacaoSegundo. A diferença é pequena: Java vence pela mão de obra e pela manutenção, e Python pelo prazo, graças ao Django. O resultado depende do empate entre prazo e mão de obra na Tabela~\ref{tab:criterios}: com o prazo um pouco mais importante (2 em vez de 1), Python passaria à frente, com 35,4\% contra 32,5\%.

\begin{table}[H]
    \centering
    \small
    \renewcommand{\arraystretch}{1.25}
    \tabelagerada{ranking}
    \caption{Ranking global: prioridade de cada alternativa em cada critério e pontuação global.}
    \label{tab:ranking}
\end{table}

\begin{thebibliography}{99}

\bibitem{arc42iso}
ARC42. \textbf{Maintainability}. Disponível em: \url{https://quality.arc42.org/qualities/maintainability}. Acesso em: 12 out. 2026. Reproduz a definição da ISO/IEC 25010:2023.

\bibitem{djangodownload}
DJANGO SOFTWARE FOUNDATION. \textbf{Download Django}. Disponível em: \url{https://www.djangoproject.com/download/}. Acesso em: 12 out. 2026.

\bibitem{djangoadmin}
DJANGO SOFTWARE FOUNDATION. \textbf{The Django admin site}. Disponível em: \url{https://docs.djangoproject.com/en/stable/ref/contrib/admin/}. Acesso em: 12 out. 2026.

\bibitem{djangoauth}
DJANGO SOFTWARE FOUNDATION. \textbf{User authentication in Django}. Disponível em: \url{https://docs.djangoproject.com/en/stable/topics/auth/}. Acesso em: 12 out. 2026.

\bibitem{adoptium}
ECLIPSE FOUNDATION. \textbf{Temurin Support}. Adoptium. Disponível em: \url{https://adoptium.net/support/}. Acesso em: 12 out. 2026.

\bibitem{gin}
GIN. \textbf{Introduction}. Gin Web Framework. Disponível em: \url{https://gin-gonic.com/en/docs/introduction/}. Acesso em: 12 out. 2026.

\bibitem{jls}
GOSLING, J. \textit{et al.} \textbf{The Java Language Specification}: Java SE 25 Edition. Cap. 4, Types, Values, and Variables. Oracle. Disponível em: \url{https://docs.oracle.com/javase/specs/jls/se25/html/jls-4.html}. Acesso em: 12 out. 2026.

\bibitem{dotnetsupport}
MICROSOFT. \textbf{.NET and .NET Core support policy}. Disponível em: \url{https://dotnet.microsoft.com/en-us/platform/support/policy/dotnet-core}. Acesso em: 12 out. 2026.

\bibitem{csharpext}
MICROSOFT. \textbf{C\#}. Extensão para o Visual Studio Code. Visual Studio Marketplace. Disponível em: \url{https://marketplace.visualstudio.com/items?itemName=ms-dotnettools.csharp}. Acesso em: 12 out. 2026.

\bibitem{csdevkit}
MICROSOFT. \textbf{C\# Dev Kit}. Visual Studio Marketplace. Disponível em: \url{https://marketplace.visualstudio.com/items?itemName=ms-dotnettools.csdevkit}. Acesso em: 12 out. 2026.

\bibitem{aspnetidentity}
MICROSOFT. \textbf{Introduction to Identity on ASP.NET Core}. Disponível em: \url{https://learn.microsoft.com/en-us/aspnet/core/security/authentication/identity}. Acesso em: 12 out. 2026.

\bibitem{efcoremigrations}
MICROSOFT. \textbf{Migrations Overview}. Entity Framework Core. Disponível em: \url{https://learn.microsoft.com/en-us/ef/core/managing-schemas/migrations/}. Acesso em: 12 out. 2026.

\bibitem{csharptypes}
MICROSOFT. \textbf{The C\# type system}. Disponível em: \url{https://learn.microsoft.com/en-us/dotnet/csharp/fundamentals/types/}. Acesso em: 12 out. 2026.

\bibitem{vscommunity}
MICROSOFT. \textbf{Visual Studio Community}. Disponível em: \url{https://visualstudio.microsoft.com/vs/community/}. Acesso em: 12 out. 2026.

\bibitem{pmi2021}
PMI - PROJECT MANAGEMENT INSTITUTE. \textbf{A Guide to the Project Management Body of Knowledge (PMBOK Guide)} and The Standard for Project Management. 7. ed. Newtown Square: Project Management Institute, 2021.

\bibitem{pmi2017}
PMI - PROJECT MANAGEMENT INSTITUTE. \textbf{Um Guia do Conhecimento em Gerenciamento de Projetos (Guia PMBOK)}. 6. ed. Newtown Square: Project Management Institute, 2017.

\bibitem{pythonversions}
PYTHON SOFTWARE FOUNDATION. \textbf{Status of Python versions}. Python Developer's Guide. Disponível em: \url{https://devguide.python.org/versions/}. Acesso em: 12 out. 2026.

\bibitem{pythontyping}
PYTHON SOFTWARE FOUNDATION. \textbf{typing: Support for type hints}. Disponível em: \url{https://docs.python.org/3/library/typing.html}. Acesso em: 12 out. 2026.

\bibitem{saaty}
SAATY, T. L. \textbf{How to Make a Decision: The Analytic Hierarchy Process}. Texto de palestra de abertura (\textit{keynote address}). [S.l.: s.n.], [s.d.]. Disponível em: \url{https://www.isahp-org.creativedecisions.net/uploads/1-how-to-make-a-decision-the-analytic-hierarchy-process-t-l-saaty.pdf}. Acesso em: 12 out. 2026.

\bibitem{springbootdb}
SPRING. \textbf{Database Initialization}. Spring Boot. Disponível em: \url{https://docs.spring.io/spring-boot/how-to/data-initialization.html}. Acesso em: 12 out. 2026.

\bibitem{springdatajpa}
SPRING. \textbf{Spring Data JPA}. Disponível em: \url{https://docs.spring.io/spring-data/jpa/reference/}. Acesso em: 12 out. 2026.

\bibitem{springsecurity}
SPRING. \textbf{Spring Security Reference}. Disponível em: \url{https://docs.spring.io/spring-security/reference/index.html}. Acesso em: 12 out. 2026.

\bibitem{so2025}
STACK OVERFLOW. \textbf{2025 Developer Survey: Technology}. Disponível em: \url{https://survey.stackoverflow.co/2025/technology/}. Acesso em: 12 out. 2026.

\bibitem{go1compat}
THE GO AUTHORS. \textbf{Go~1 and the Future of Go Programs}. Disponível em: \url{https://go.dev/doc/go1compat}. Acesso em: 12 out. 2026.

\bibitem{gorelease}
THE GO AUTHORS. \textbf{Release History}. Disponível em: \url{https://go.dev/doc/devel/release}. Acesso em: 12 out. 2026.

\bibitem{gospec}
THE GO AUTHORS. \textbf{The Go Programming Language Specification}. Disponível em: \url{https://go.dev/ref/spec}. Acesso em: 12 out. 2026.

\bibitem{ach2003}
UNIVERSIDADE DE SÃO PAULO. \textbf{ACH2003 - Computação Orientada a Objetos}. JupiterWeb. Disponível em: \url{https://uspdigital.usp.br/jupiterweb/obterDisciplina?sgldis=ACH2003}. Acesso em: 12 out. 2026.

\bibitem{gradesi}
UNIVERSIDADE DE SÃO PAULO. \textbf{Grade curricular: Bacharelado em Sistemas de Informação}. JupiterWeb. Currículo com início em 1º jan. 2027. Disponível em: \url{https://uspdigital.usp.br/jupiterweb/listarGradeCurricular?codcg=86&codcur=86200&codhab=202&tipo=N}. Acesso em: 12 out. 2026.

\bibitem{uspdiretrizes}
UNIVERSIDADE DE SÃO PAULO. Coordenadoria de Administração Geral. \textbf{Diretrizes Orçamentárias e Revisão do Planejamento Plurianual 2026}. São Paulo: USP, 2025. Disponível em: \url{https://sites.usp.br/codage/wp-content/uploads/sites/264/2026/01/2026_Diretrizes-Orcamentarias.pdf}. Acesso em: 12 out. 2026.

\bibitem{vargas2010}
VARGAS, R. V. Using the Analytic Hierarchy Process (AHP) to Select and Prioritize Projects in a Portfolio. In: PMI GLOBAL CONGRESS 2010 - NORTH AMERICA, Washington, DC. \textbf{Anais}. Newtown Square: Project Management Institute, 2010. Disponível em: \url{https://ricardo-vargas.com/articles/analytic-hierarchy-process/}. Acesso em: 12 out. 2026.

\end{thebibliography}

\end{document}
